% ============================================================
%  第1章：绪论
% ============================================================
\chapter{绪论}

\section{研究背景与意义}

\zhlipsum[1]

近年来，Vaswani等人\cite{vaswani2017attention}提出的 Transformer 架构完全摒弃了循环结构，仅依赖自注意力机制实现序列建模，在机器翻译任务上取得了当时最优的结果。这一突破引发了自然语言处理领域的范式转变\cite{devlin2019bert}。

\zhlipsum[2]

\section{国内外研究现状}

\zhlipsum[3]

在此基础上，Devlin等人\cite{devlin2019bert}进一步提出了 BERT 预训练模型，通过在海量无标注语料上预训练深度双向 Transformer，再在下游任务上微调，在十一项自然语言处理基准上全面刷新了记录。多项后续工作\cite{vaswani2017attention,devlin2019bert}均沿用了这一预训练-微调范式。

\zhlipsum[4]

值得注意的是，原始 Transformer 的自注意力复杂度为 $O(n^2 \cdot d)$\cite{vaswani2017attention}，这限制了其在长序列场景中的应用。

\section{研究内容与创新点}

\zhlipsum[6]

针对上述问题，本文的主要贡献包括：

\begin{enumerate}
  \item 提出稀疏注意力与局部敏感哈希相结合的混合编码策略，在保持模型表达能力的同时显著降低自注意力计算复杂度。
  \item 设计面向长序列的高效训练方案，通过梯度累积与混合精度训练减少显存占用。
  \item 在机器翻译和文本摘要等多个基准上系统验证所提方法的有效性。
\end{enumerate}

\section{论文组织结构}

\zhlipsum[11]
